\chapter{Inhibitor Destruction by Moist Conditions}

\textbf{Destruction of the inhibiting system takes place at 40.} In many
species the germination inhibitors are destroyed simply by exposure to moisture.
There is an induction period before germination begins. After the induction period
there is a sudden onset of germination. This pattern can occur at either 40 or 70.
Some examples of this behavior for species germinating at 40 are shown in Table 6-1.
This group is a problem only because it has not been generally recognized that many
species germinate exclusively at 40.

A pattern that requires somewhat more planning to achieve germination are
those in which the inhibiting system is destroyed by exposure to three months moist at
40 followed by germination on a shift to 70. These will be termed 40-70 germinators.
Some examples of these are given in Table 6-2. This pattern is well known, but the
period at 40 has traditionally been spoken of as a period required to "break
dormancy." It should be clear from the discussion in Chapters 1 and 2 that this is a
\textit{conditioning period} and a period of \textit{maximum metabolic activity} for those chemical
processes that destroy the inibiting system. It is only a dormant period in respect to
vegetative growth.

The destructions of the inhibitor systems at 40 are rate processes and will be
different for each species. Minimum times for the low temperature cycle for seeds
have been reported to be twelve hours for wheat, seven days for Poa annua, fourteen
days for Delphinium ambiguum, and six months for Crataegus mollis (ref. 13). Data
given below in Table 6-3 for Nemastylis acuta shows that the time for this species is
about three months.

The destructions of germination inhibitors in seeds at 40 are clearly related to.
the destruction of growth inhibitors which occurs during winter in \textit{vegetatively dormant}
flower buds and vegetative buds of deciduous trees and shrubs as well as \textit{vegetatively
dormant} bulbs and tubers. A similar period of about three months at temperatures
around 35-42 degrees F is needed in order for vegetative growth to occur on shifting to
warmer temperatures. The similarity or possible identity with the chemical processes
in seeds remains for future investigators to determine. In any event such periods
should be called conditioning periods and not breaking dormancy unless one
specifies that one is speaking of dormancy only in regard to vegetative growth.
The fact that these inhibitor destruction reactions occur at 40 but not 70 requires
that the rates of the processes increase enormously in lowering the temperature from
70 to 40. Enzyme reactions and protein and membrane restructuring have a special
capacity for these unusual negative temperature coefficients, and the theory of this is
explained in Chapter 19.

The data in Tables 6-1 and 6-2 represent two limiting situations. The inhibitors
are destroyed at 40 in both, but in the examples in Table 6-1 germination occurs at 40
whereas in the examples in Table 6-2 germination is so slow at 40 relative to the
germination rate at 70 that a shift to 70 is necessary to achieve germination.
There are intermediate situations where the rate of germination at 40 is
significant, but the rate at 70 is much faster. Nemastylis acuta illustrates this situation
and the data appear in Table 6-3. The inhibitor destruction requires about three
months at 40. At this time there is the option for either allowing a slower germination to
proceed at 40 or a faster germination to proceed at 70 by shifting the seed to 70.
Analysis of the data shows that the germination rate at 70 is 2.3 times faster than at 40.
This temperature coefficient is in the range typical for most chemical processes. Note
that the grower will fail if the seeds are kept moist at 40 in the dark, because after three
months they will start to germinate and die for lack of light. Incidentally, Nemastylis
acuta has beautiful two inch wide sky blue flowers with Tigridialike pleated foliage. It
is one of the most beautiful of hardy bulbs and is little known in horticulture.

\textbf{Destruction of the inhibiting system takes place at 70.} This situation
also generates two limiting germination patterns. In one pattern there is a long
induction period at 70 followed ultimately by germination at 70. The interpretation is
that inhibitors are destroyed at -70 and that the germination itself is faster at 70. In the
other limiting pattern the germination inhibitors are destroyed at 70, but germination
occurs only at 40 so that after the inhibitors have been destroyed at 70, the seeds must
be shifted to 40 in order for germination to occur. Species with this pattern are called
70-40 germinators. This type of pattern has not been explicitly recognized in the
literature, and it was a surprise to find that there were comparable numbers of 70-40
and 40-70 germinators.

Following is a list of species in which germination occurs at 70, but only after an
induction period. This list is restricted to species where the induction time at 70 was
two weeks or longer since this length of time makes it reasonably certain that
germination inhibitors were present. The induction is given in weeks (w) in
parentheses after the name of each species. Aucuba japonica (5 w), Centaurium
erythrosa (2.5 w), Clematis integrifolia (6 w), Clematis pitcheri (4 w), Clematis recta (8
w), Clematis viorna (8 w), Cortusa matthioli (3 w), Cortusa turkestanica (3 w),
Delosperma cooperi (4 w), Eritrichum howardi (4 w), Haberlea rhodopensis (4 w),
Jankae heldrichii (4 w), Leucojum vemum (8 w), Lilium auratum (3 w), Lilium
canadense editorum (7 w), Lilium canadense flavum (7 w), Lilium cernuum (6 w),
Lilium concolor (8 w), Lilium davidi (8 w), Lilium maximowiczii (8 w), Paeonea
officinalis (12 w), Paeonea suffruticosa (12 w), Podophyllum emodi (5 w), Prunus
armenaica (5 w), Silene polypetala (12 w), Silene regia (12 w), Thlaspi bulbosum
(6 w), Thlaspi montanum (6 w), and Thlaspi rotundifolia (6 w).

Table 6-4 lists those species that were straightforward 70-40 germinators.
Examples that had additional complexities are not listed here, but are of course
described in Chapter 20. Note that usually the germination was completed rapidly
over a time span of one to three weeks and this was shorter than the induction times.
There was a minor difference in the species that germinated at 70 relative to the
species that germinated at 40. Where germination occurred at 70, development of
photosynthesizing cotyledons or leaves soon followed. Where germination occurred
at 40, sometimes cotyledons or leaves developed at 40 as in Brodiae congesta and
Eranthis hyemalis, but more commonly only the radicle developed at 40, and a shift to
70 was required for leaf development.

In rare examples either dry storage of the seed for six months at 70 or
subjecting the seed to three months moist at 40 were equally effective in destroying
the germination inhibitors. In Eupatorium urticaefolium germination at 70 was 95\%
complete in 2-5 days after the three months at 40 and 95\% complete in 5-10 days after
the dry storage. The slightly longer times in the latter were due to time required for
water transport in the freshly moistened seed.

Up to now it has been emphasized that inhibitor destruction under moist
conditions virtually never occurred at both 40 and 70. Anemone baldensis is
possible exception in that germination occurred at both 40 and 70 and with significant
induction times. Seed sown at 40 had an induction time of 44 days and a rate of
germination of 4.0\%/d. Seed sown at 70 had an induction time of 18 days and a
germination rate of 5.1\%/d. These values are for fresh seed. As with Eupatorium
urticaefolium, dry storage removed most of the inhibition so that germination at 70 was
speeded up and was 95\% complete in 10-15 days.

A major difference was present between the 40-70 germinators and the 70-40
germinators. In the 40-70 germinators germination took place immediately after the
shift to 70 so that there was usually no significant induction period at 70. In the 70-40
germinators germination took place at 40 usually after a significant induction period,
and this induction period was followed by relatively rapid germination over a period of
two to three weeks. This is interpreted to mean that there was a second set of
germination inhibitors that had to be destroyed at 40. This leads us into the concept of
two or more sets of germination inhibitors which is the subject of the next chapter.

\begin{table}[h!]
\refstepcounter{table}

\begin{center}
Table \thetable. Inhibitor Destruction at 40 Followed by Germination at 40 (a,b)
\end{center}

\begin{tabular}{|l|l|l|}
\hline
Species & Induction time & Germination rate \\
\hline
Allium albopilosum & 5 weeks & 85\% in 5-13 weeks \\
Allium moly & 9 weeks & 90\% in 9-11 weeks \\
Calandrinia caespitosa & 7 weeks & 40(33\%)-70(29\% in 6-10 days)(c) \\
Camassia leichtlinii & 5 weeks & 100\% in 5-14 weeks \\
Chaenomeles japonica & 7 weeks & 86\%, \%5/d \\
Disporum lanuginosum & 7 weeks & 100\% in 8th week \\
Erythronium citrinum & 10 weeks & 40(65\%)-70(35\% in 1-2 days)(c) \\
Erythronium hendersoni & 8 weeks & 93\% in 8-12 weeks \\
Euonymus bungeana & 5 weeks & 90\% in 6th week \\
Fritillaria acmopetala & 9 weeks & 40(45\%)-70(23\% in 2-6 days)(c) \\
Fritillaria persica & 8 weeks & 97\%, ind. 54 days, 8\%/d \\
Gladiolus anatolicus & 7 weeks & 90\% in 7-9 weeks \\
Gladiolus imbricata & 7 weeks & 90\% in 7-9 weeks \\
Gladiolus kotschyanus & 6 weeks & 90\% in 6-12 weeks \\
Hebe guthreana & 7 weeks & 96\%, ind. t 48 days, 4\%/d \\
Iris bracteata & 8 weeks & 90\% in 9th week \\
Iris germanica & 7 weeks & 40(31\%)-70(31\% in 1-6 days)(c) \\
Iris magnifica & 10 weeks & 30-100\% in 10-12 weeks \\
Kirengoshima palmata & 12 weeks & 90\% in 12-16 weeks \\
Lonicera sempervirens & 7 weeks & 100\%, ind. t 50 days, 3\%/d \\
Parrotia persica & 11 weeks & 90\% in 12th week \\
Penstemon glabrescens & 5 weeks & 84\% in 5-12 weeks \\
Rosa multiflora & 5 weeks & 26\% in 5-7 weeks \\
Saponaria caespitosa & 9 weeks & 87\% in 9-12 weeks (dry stored seed) \\
Scilla hispanica & 9 weeks & 32\% in 10th week \\
Zigadenus fremonti & 6 weeks & 92\% in 6-9 weeks \\
\hline
\end{tabular}

(a) It is important to remember that these species gave no germination when
sown at 70 and kept at 70. Where the percent germination was low, it was because of
a predominance of empty seed coats and not because there was further germination
at a later date.

(b) The Table includes only those species where the induction time was five
weeks or longer since it was felt that these were the clearest examples of an induction
period.

(c) Probably germination would have been completed at 40 if more time had
been given. These are examples where the germination rate is much faster at 70 than
at 40.

\end{table}

\begin{table}[h!]
\refstepcounter{table}
\begin{center}
Table \thetable. Inhibitor Destruction at 40 Followed by Rapid Germination on Shifting to 70
\end{center}
\par

\begin{tabular}{|l|l|}
\hline
Species & Germination rate after the shift to 70 \\
\hline
Aesculus hippocasteanum (a) & 100\% in 1-7 days \\
Aesculus pavia & 100\% in 1-7 days \\
Allium tricoccum & 95\%, largely in 10-13 weeks \\
Aquilegia formosa & 50\% in 14-42 days \\
Asciepias tuberosa & 100\% in 2-7 days \\
Asimina triloba & 76\% in 2-35 days \\
Aureolaria virginica & 30\% in 5-24 days \\
Campanula altaica & 68\% in 4-14 days \\
Ceanothus americana & 28\% in 1-21 days \\
Cercidiphyllum japonicum & 50\%, ind. t 2 days, 10\%/d \\
Cornus nuttailli & 40\% in 8-10 days \\
Diospyros virginiana & 60\% in 7-28 days \\
Dodecatheon jeffreyi & 90\% in 1-5 days \\
Dodecatheon media & 92\% in 1-3 days \\
Eupatorium purpureum & 90\% in 3-5 days \\
Ilex japonica & 87\% in 90\% in 1-2 days \\
Myrica pennsylvanica & 85\% in 1-2 days if wax was removed from seed \\
Phytolacca americana & 80\% in 5-16 days if seed was washed and dry stored \\
Primula intercedens & 43\% in 6-10 days if seed was dry stored \\
Primula japonica & 100\% in 2-4 days if seed was dry stored one month \\
Primula parryi & 32\% in 4-42 days \\
Rhamnus caroliniana & 98\% in 4-7 days \\
Ruellia humihis & 33\% in 3-9 days (rest of the seed rots quickly) \\
Trillium albidum & 100\% in 2-4 days \\
Trohlius acaulis & 70\% in 6-10 days \\
Vernonia novaboracensis & 100\% in 2-4 days \\
Viburnum tomentosum & 50\% in 1-26 days \\
Vitis vulpina & 76\% in 4-6 days \\
\hline
\end{tabular}

\hfil

(a) With large fleshy seeds like this species, there is apparently sufficient
moisture in the seed itself so that dry storage at 40 destroys the inhibitors as effectively
as 3 m moist at 40. Although 6 m dry storage at 40 was used, it is likely that 3 m would.
have been equally effective.

\label{table:id40}
\end{table}

\begin{table}[h]
\refstepcounter{table}
\begin{center}
Table \thetable. Inhibitor Destruction at 40 Followed by Germination at Either 40 or 70 for
Nemastylis Acuta
\end{center}
\par
\begin{tabular}{|l l|l|l|l|l|l|}
\hline
Time in months (at 40)        & 1 & 2 & 3    & 4    & 4.5     & 5 \\
\hline
Percent germination           & & & & & & \\
\quad seed kept at 40         & 0 & 0 & 0    & 38\% & 58\%    & 72\% \\
\quad seed shifted, to 70 (a) & 0 & 0 & 44\% & 92\% & 90\%(b) & 90\%(b) \\
\quad seed kept at 70         & 0 & 0 & 0    & 0    & 0       &  0 \\
\hline
\end{tabular}

\hfil

(a) The shift to 70 was made after the time at the head of the column. The
germination at 70 was rapid and was completed in 2-7 days. The value of 44\% at the
end of the three months at 40 is very sensitive to the exact conditions, and values
ranging from 40-80\% were observed. This is because the inhibitor destruction
reaction is just at the point of being completed.

(b) These values include the seeds that had germinated earlier at 40.

\label{table:acuta}
\end{table}

\begin{table}[h]
\refstepcounter{table}
\begin{center}
Table \thetable. Inhibitor Destruction at 70 Followed by Germination at 40 (a)
\end{center}

\par
\begin{tabular}{|l|l|}
\hline
Species             & Germination rate at 40 after 3 months at 70 \\
\hline
Anemone blanda      & 75\%, ind. t 63 d, 3\%/d \\
Actaea pachypoda      & 85\% in 5-8 w       \\
Actaea spicata      & 88\% in 4-7 w         \\
Anemone blanda      & 75\%, ind. t 63 d, 3d \\
Chionodoxa luciliae & 50\% in 12th w        \\
Cimicifuga: racemosa & 100\% in 3-6w        \\
Claytonia virginica &  79\% in 7th w        \\
Clematis alpina     & 80\% in 5-11 w        \\
Clematis maculata    & 100\% in 4-8w        \\
Colchicum autumnale  & 33\% in 6th w        \\
Colchicum luteum     & 58\% in 0-3 w        \\
Corydalis lutea     & 70\% in 10-12 w       \\
Crocus speciosus     & 86\% in 3-5 w        \\
Delphinium glaucum     & 25\% in 4-6w       \\
Eranthishyemalis       & 100\% in 8-9w      \\
Erythronium americanum & 60\% in 10-12w (b) \\
Erythronium revolutum  & 33\% in 12th w     \\
Fritillaria thunbergii & 50\% in 5-7w \\
Galanthus elwesii & 45\% in 5th w \\
Galanthus nivalis & 45\% in 2-5 w \\
Helleborus niger & 93\% in 9th w \\
Helleborus orientalis & 100\% in 5-10 w \\
Hyacinthus orientalis & 100\% in 6-8w \\
Iris albomarginata & 50\% in 7-12 w \\
Isopyrum biternatum & 5\% in 8-11 w \\
Ligustrum abtusifolium & 4\% in 2-12 w \\
Lilium pardalinum & 0\% in 2nd w \\
Lilium parvurn & 90\% in 2nd w \\
Lilium shastense & 90\% in 2nd w \\
Lilium washingtonianum & 61\% in 3-12 w \\
Melampyrum Iineare & 55\% in 4-8w \\
Scilla campanulata & 75\% in 3-11 w \\
Scilla tubergeniana & 100\% in 3-5 w \\
\hline
\end{tabular}
\par
(a) No further germination occurred unless noted.

(b) Germination was 100\% after six cycles.
\label{table:g7040}
\end{table}

